%============================================================
% CH2 · P5 · PROBLEM 1 — ALL ANCESTORS OF A NODE IN A DAG (LeetCode 2192)
%============================================================
\subsection{Problem 2.5.1: All Ancestors of a Node in a Directed Acyclic Graph}
\problemlabel{2.5.1}{ssec:ch2-all-ancestors}

% the sample DAG of this problem as a pic; vertex styles g2p5-a<k>, edge styles g2p5-f<uv>
\tikzset{
    g2p5-a0/.style=gv, g2p5-a1/.style=gv, g2p5-a2/.style=gv, g2p5-a3/.style=gv,
    g2p5-a4/.style=gv, g2p5-a5/.style=gv, g2p5-a6/.style=gv, g2p5-a7/.style=gv,
    g2p5-f03/.style=gde, g2p5-f04/.style=gde, g2p5-f13/.style=gde, g2p5-f24/.style=gde, g2p5-f27/.style=gde,
    g2p5-f35/.style=gde, g2p5-f36/.style=gde, g2p5-f37/.style=gde, g2p5-f46/.style=gde,
    g2p5-allrev/.style={g2p5-f03/.style=g2p5-re, g2p5-f04/.style=g2p5-re, g2p5-f13/.style=g2p5-re,
        g2p5-f24/.style=g2p5-re, g2p5-f27/.style=g2p5-re, g2p5-f35/.style=g2p5-re, g2p5-f36/.style=g2p5-re,
        g2p5-f37/.style=g2p5-re, g2p5-f46/.style=g2p5-re},
    g2p5-S1/.pic={
        \node[g2p5-a1] (s1) at (0,2) {1}; \node[g2p5-a0] (s0) at (1.6,2) {0}; \node[g2p5-a2] (s2) at (3.2,2) {2};
        \node[g2p5-a3] (s3) at (0.8,1) {3}; \node[g2p5-a4] (s4) at (2.4,1) {4};
        \node[g2p5-a5] (s5) at (0,0) {5}; \node[g2p5-a6] (s6) at (1.6,0) {6}; \node[g2p5-a7] (s7) at (3.2,0) {7};
        \draw[g2p5-f03] (s0) -- (s3); \draw[g2p5-f04] (s0) -- (s4); \draw[g2p5-f13] (s1) -- (s3);
        \draw[g2p5-f24] (s2) -- (s4); \draw[g2p5-f27] (s2) -- (s7); \draw[g2p5-f35] (s3) -- (s5);
        \draw[g2p5-f36] (s3) -- (s6); \draw[g2p5-f37] (s3) -- (s7); \draw[g2p5-f46] (s4) -- (s6);
    }
}

\problemcard{2.5.1}{All Ancestors of a Node in a Directed Acyclic Graph}{LeetCode 2192}{https://leetcode.com/problems/all-ancestors-of-a-node-in-a-directed-acyclic-graph/}{\exmedium\ \lc{Medium}}{For every vertex, one search in the reversed graph collects its ancestors; scanning the marks in index order gives sorted lists for free}

\subsubsection*{1. Problem statement}

A directed acyclic graph has $n$ vertices $0, \dots, n-1$ and a list of edges $[\mathit{from}, \mathit{to}]$. A vertex $u$ is an \textbf{ancestor} of $v$ if $u$ can reach $v$ along the edges. Return, for every $v$, the list of its ancestors in increasing order.

\textbf{Constraints:} $1 \le n \le 1000$, $0 \le m \le \min(2000, n(n-1)/2)$; no self-loops, no duplicate edges; the graph is acyclic.

\begin{table}[H]
\centering\small
\begin{tabular}{@{}p{5.6cm}p{5.6cm}l@{}}
\toprule
\textbf{Input} & \textbf{Official output} & \textbf{Ours} \\
\midrule
$n = 8$, \texttt{[[0,3],[0,4],[1,3],[2,4],}\newline\texttt{[2,7],[3,5],[3,6],[3,7],[4,6]]} & \texttt{[[],[],[],[0,1],[0,2],[0,1,3],}\newline\texttt{[0,1,2,3,4],[0,1,2,3]]} & identical \\
$n = 5$, all 10 edges $i \to j$ with $i < j$ & \texttt{[[],[0],[0,1],[0,1,2],}\newline\texttt{[0,1,2,3]]} & identical \\
\bottomrule
\end{tabular}
\end{table}

\begin{diagrambox}
\begin{tikzpicture}[x=1.15cm, y=1.1cm]
    \begin{scope}
        \pic {g2p5-S1};
        \node[font=\scriptsize\bfseries] at (1.6,2.75) {Sample 1};
    \end{scope}
    \begin{scope}[xshift=5.6cm, g2p5-a6/.style=gvdst, g2p5-a0/.style=gvsel, g2p5-a1/.style=gvsel,
                  g2p5-a2/.style=gvsel, g2p5-a3/.style=gvsel, g2p5-a4/.style=gvsel,
                  g2p5-f03/.style=gpath, g2p5-f13/.style=gpath, g2p5-f36/.style=gpath,
                  g2p5-f04/.style=gpath, g2p5-f24/.style=gpath, g2p5-f46/.style=gpath]
        \pic {g2p5-S1};
        \node[font=\scriptsize\bfseries] at (1.6,2.75) {the ancestors of 6};
    \end{scope}
    \begin{scope}[xshift=11.2cm, x=0.95cm, y=0.7cm]
        \foreach \v/\a [count=\i from 0] in {0/{---}, 1/{---}, 2/{---}, 3/{0 1}, 4/{0 2}, 5/{0 1 3}, 6/{0 1 2 3 4}, 7/{0 1 2 3}}
            \node[font=\scriptsize, anchor=west] at (0,2.4-0.59*\i) {\textbf{\v:}\ \a};
        \node[font=\scriptsize\bfseries] at (1.1,3.05) {answer};
    \end{scope}
\end{tikzpicture}
\figcap{Sample 1. Vertex 6 is reached from 3 and from 4, hence from all of 0, 1, 2, 3, 4 (purple, orange paths). Vertices 0, 1, 2 have in-degree 0 and no ancestors.}
\end{diagrambox}

\subsubsection*{2. The graph model}

\begin{modelbox}[Graph Model --- All Ancestors]
\renewcommand{\arraystretch}{1.35}
\begin{tabularx}{\linewidth}{@{}>{\bfseries}l X@{}}
Vertices & the $n$ given vertices \\
Edges & directed, as given; the graph is a DAG (we will not need that) \\
Weights & none \\
Question & for every $v$: the set $\{u \ne v : u \leadsto v\}$ $=$ the vertices reachable from $v$ in the \textbf{reversed} graph $G^R$, sorted \\
Size & $n \le 1000$, $m \le 2000$; the output alone can have $\binom{n}{2} \approx 5 \cdot 10^5$ numbers, so \bigO{n(n+m)} $\approx 3 \cdot 10^6$ is the natural target \\
\end{tabularx}
\end{modelbox}

\begin{diagrambox}
\begin{tikzpicture}[x=1.15cm, y=1.1cm]
    \begin{scope}
        \pic {g2p5-S1};
        \node[font=\scriptsize\bfseries] at (1.6,2.75) {$G$: ``who reaches $v$?''};
    \end{scope}
    \draw[flow] (4.0,1) -- node[lbl, above]{turn every edge} (5.6,1);
    \node[note, text width=2.2cm, align=center] at (4.8,0.25) {\texttt{radj[to]}\\\texttt{.push\_back(from)}};
    \begin{scope}[xshift=7.0cm, g2p5-allrev]
        \pic {g2p5-S1};
        \node[font=\scriptsize\bfseries] at (1.6,2.75) {$G^R$: ``whom does $v$ reach?''};
    \end{scope}
\end{tikzpicture}
\figcap{The model: ancestors in $G$ are descendants in $G^R$. A plain reachability search from $v$ in $G^R$ is all we need.}
\end{diagrambox}

\subsubsection*{3. Key observations}

\begin{itemize}[leftmargin=*]
    \item \textbf{Ancestors of $v$ $=$ vertices reachable from $v$ in $G^R$} (the lemma of the pattern). One search per vertex: $n$ searches of \bigO{n + m} each.
    \item \textbf{The output is quadratic anyway.} A chain $0 \to 1 \to \dots \to 999$ has $\sum_v v = 499\,500$ ancestor entries; nothing can be faster than writing them. So \bigO{n(n+m)} is not ``slow'', it is the size of the answer plus a factor $\le 3$.
    \item \textbf{Sorted for free.} After the search from $v$, scan $u = 0, 1, \dots, n-1$ and collect the marked ones: the list is sorted without \texttt{sort}. Collecting in \emph{visiting} order would not be: from 6 the search visits $3, 4, 0, 2, 1$.
    \item \textbf{One marker array, $n$ searches.} Instead of clearing \texttt{seen} before every search, store \texttt{mark[u] = v} (``seen in the search from $v$''): no \bigO{n} reset, no reallocation.
    \item \textbf{Acyclicity is not needed} for the search approaches. It only matters for the topological-order alternative.
\end{itemize}

\subsubsection*{4. Approaches}

\begin{bruteforcebox}
\textbf{Idea: ask every pair.} For every $v$ and every $u \ne v$, run a fresh search from $u$ and check whether it reaches $v$.

\begin{lstlisting}[style=cpp]
#include <bits/stdc++.h>
using namespace std;

class Solution {
    vector<vector<int>> adj;
    bool reach(int s, int t) {                    // a fresh iterative DFS from s
        vector<char> seen(adj.size(), 0);
        vector<int> st(1, s);
        seen[s] = 1;
        while (!st.empty()) {
            int u = st.back(); st.pop_back();
            if (u == t) return true;
            for (int v : adj[u]) if (!seen[v]) { seen[v] = 1; st.push_back(v); }
        }
        return false;
    }
public:
    vector<vector<int>> getAncestors(int n, vector<vector<int>>& edges) {
        adj.assign(n, vector<int>());
        for (auto& e : edges) adj[e[0]].push_back(e[1]);
        vector<vector<int>> ans(n);
        for (int v = 0; v < n; v++)
            for (int u = 0; u < n; u++)                 // increasing u: already sorted
                if (u != v && reach(u, v)) ans[v].push_back(u);
        return ans;
    }
};
\end{lstlisting}

\begin{complexitybox}
\textbf{Time:} $n^2$ searches of \bigO{n+m}: \bigO{n^2(n+m)}. On a chain of $1000$ vertices plus $1001$ random forward edges ($m = 2000$) it took \textbf{5.0\,s} on our machine. \textbf{Space:} \bigO{n+m}. \textbf{Verdict:} \tle; kept as the reference for testing.
\end{complexitybox}
\end{bruteforcebox}

\begin{algorithmbox}[Approach 2: one reverse search per vertex]
\textbf{Idea.} Build \texttt{radj}. For every $v$, search from $v$ in $G^R$, marking with \texttt{mark[u] = v}; then scan $u$ in increasing order and collect the marked $u \ne v$.

\begin{lstlisting}[style=cpp]
#include <bits/stdc++.h>
using namespace std;

class Solution {
public:
    vector<vector<int>> getAncestors(int n, vector<vector<int>>& edges) {
        vector<vector<int>> radj(n);
        for (auto& e : edges) radj[e[1]].push_back(e[0]);   // from -> to, kept at "to"
        vector<vector<int>> ans(n);
        vector<int> mark(n, -1), st;                      // mark[u] = v: seen from v
        for (int v = 0; v < n; v++) {
            mark[v] = v;
            st.assign(1, v);
            while (!st.empty()) {                               // iterative DFS in G^R
                int x = st.back(); st.pop_back();
                for (int u : radj[x])
                    if (mark[u] != v) { mark[u] = v; st.push_back(u); }
            }
            for (int u = 0; u < n; u++)                   // index order = sorted
                if (u != v && mark[u] == v) ans[v].push_back(u);
        }
        return ans;
    }
};
\end{lstlisting}

\begin{complexitybox}
\textbf{Time:} $n$ searches of \bigO{n+m} plus $n$ scans of \bigO{n}: \bigO{n(n+m)} $\approx 3 \cdot 10^6$. The same chain input: \textbf{13\,ms}. \textbf{Space:} \bigO{n+m} besides the output. \textbf{Verdict:} fast enough by a factor of several hundred. The search is iterative, so even a chain of $10^5$ vertices would not overflow the stack.
\end{complexitybox}
\end{algorithmbox}

\begin{invariantbox}[Why it is correct]
\textbf{Claim.} After the search from $v$, $\texttt{mark}[u] = v$ iff $u \leadsto v$ in $G$.

\emph{Only ancestors are marked.} $v$ is marked first. A vertex $u$ is marked only from a popped $x$ with $u \in \texttt{radj}[x]$, i.e.\ an edge $u \to x$; by induction $x \leadsto v$, so $u \to x \leadsto v$.

\emph{Every ancestor is marked.} Suppose $u \leadsto v$ along $u = y_0 \to y_1 \to \dots \to y_k = v$ but $u$ is unmarked. Take the largest $i$ with $y_i$ unmarked; $y_{i+1}$ is marked, was pushed and later popped, and then it scanned $\texttt{radj}[y_{i+1}] \ni y_i$ and would have marked it. Contradiction.

\emph{Old marks do not interfere.} Marks from earlier searches hold values $< v$, so ``\texttt{mark[u] != v}'' is exactly ``not yet seen in this search''. The final scan in index order outputs the set sorted. \hfill$\blacksquare$
\end{invariantbox}

\begin{alternativebox}[Approach 3: forward searches that hand out their name]
\textbf{Idea.} Turn it around once more: run a forward search from every $u = 0, 1, \dots, n-1$ \emph{in this order}, and append $u$ to the list of every vertex it reaches. Since $u$ increases, every list receives its entries in increasing order --- sorted for free again, no reversed graph needed. Recursion depth is at most $n \le 1000$, so the recursive form is safe here.

\begin{lstlisting}[style=cpp]
#include <bits/stdc++.h>
using namespace std;

class Solution {
    vector<vector<int>> adj, ans;
    vector<int> mark;                                  // mark[x] = u: x reached from u
    void dfs(int x, int u) {
        for (int v : adj[x])
            if (mark[v] != u) { mark[v] = u; ans[v].push_back(u); dfs(v, u); }
    }
public:
    vector<vector<int>> getAncestors(int n, vector<vector<int>>& edges) {
        adj.assign(n, vector<int>());
        ans.assign(n, vector<int>());
        mark.assign(n, -1);
        for (auto& e : edges) adj[e[0]].push_back(e[1]);
        for (int u = 0; u < n; u++) { mark[u] = u; dfs(u, u); }   // u increasing
        return ans;
    }
};
\end{lstlisting}

\begin{center}
\begin{tikzpicture}[x=1.15cm, y=1.1cm]
    \begin{scope}[g2p5-a0/.style=gvsrc, g2p5-a3/.style=gvdone, g2p5-a5/.style=gvdone, g2p5-a6/.style=gvdone,
                  g2p5-a7/.style=gvdone, g2p5-a4/.style=gvdone,
                  g2p5-f03/.style={gtree, -{Stealth[length=2.4mm]}}, g2p5-f35/.style={gtree, -{Stealth[length=2.4mm]}},
                  g2p5-f36/.style={gtree, -{Stealth[length=2.4mm]}}, g2p5-f37/.style={gtree, -{Stealth[length=2.4mm]}},
                  g2p5-f04/.style={gtree, -{Stealth[length=2.4mm]}}, g2p5-f46/.style=gfaded]
        \pic {g2p5-S1};
        \node[font=\scriptsize\bfseries] at (1.6,2.75) {forward search from $u = 0$};
    \end{scope}
    \node[note, text width=7.2cm, align=left, anchor=west] at (4.3,1.0) {visits $3, 5, 6, 7, 4$ and appends 0 to each of their lists; the edge $4 \to 6$ (grey) finds 6 already marked with 0, so 0 is \emph{not} appended twice.\\[3pt]
    Visiting orders of the next searches: $u=1$: $3,5,6,7$; $u=2$: $4,6,7$; $u=3$: $5,6,7$; $u=4$: $6$; $u = 5, 6, 7$: nothing.\\[3pt]
    List of 6 grows as $[0] \to [0,1] \to [0,1,2] \to [0,1,2,3] \to [0,1,2,3,4]$.};
\end{tikzpicture}
\end{center}

\begin{complexitybox}
\textbf{Time:} \bigO{n(n+m)}, 15\,ms on the chain input. \textbf{Space:} \bigO{n+m} plus recursion depth $\le n$. \textbf{Verdict:} as good as Approach~2. Which one to write is a matter of taste: Approach~2 needs \texttt{radj} and a scan, Approach~3 needs the increasing order of $u$.
\end{complexitybox}
\end{alternativebox}

\begin{tricksbox}[Approach 4: Topological order and bitsets]
In a DAG, ancestors can be \emph{merged} instead of searched: if $u \to v$, then every ancestor of $u$ and $u$ itself are ancestors of $v$. Processing the vertices in topological order (Kahn's algorithm, Chapter~5) guarantees that $\mathrm{anc}(u)$ is complete before it is merged into its successors. With \texttt{bitset<1000>} a merge is $1000/64 \approx 16$ word operations.

\begin{lstlisting}[style=cpp]
#include <bits/stdc++.h>
using namespace std;

class Solution {
public:
    vector<vector<int>> getAncestors(int n, vector<vector<int>>& edges) {
        vector<vector<int>> adj(n);
        vector<int> indeg(n, 0), order;
        for (auto& e : edges) { adj[e[0]].push_back(e[1]); indeg[e[1]]++; }
        for (int v = 0; v < n; v++) if (!indeg[v]) order.push_back(v);
        for (size_t i = 0; i < order.size(); i++)         // Kahn: topological order
            for (int v : adj[order[i]]) if (--indeg[v] == 0) order.push_back(v);
        vector<bitset<1000>> A(n);                        // A[v][u]: u is an ancestor
        for (int u : order)
            for (int v : adj[u]) { A[v] |= A[u]; A[v][u] = 1; }   // anc(u) + u into v
        vector<vector<int>> ans(n);
        for (int v = 0; v < n; v++)
            for (int u = 0; u < n; u++) if (A[v][u]) ans[v].push_back(u);
        return ans;
    }
};
\end{lstlisting}
\textbf{Time:} \bigO{n + m \cdot n/64 + n^2} --- the final $n^2$ scan dominates; 7\,ms on the chain input. \textbf{Space:} $n^2$ bits $= 125$\,KB. Unlike the searches it needs the graph to be acyclic, and a bitset size fixed at compile time; for $n = 10^5$ the $n^2$ bits (1.25\,GB) would not fit.
\end{tricksbox}

\subsubsection*{5. Trace}

\begin{tracebox}[{Sample 1, Approach 2: the searches from $v = 6$ and $v = 7$}]
\texttt{radj} in input order: $\texttt{radj}[3] = 0, 1$;\ $\texttt{radj}[4] = 0, 2$;\ $\texttt{radj}[5] = 3$;\ $\texttt{radj}[6] = 3, 4$;\ $\texttt{radj}[7] = 2, 3$; the lists of $0, 1, 2$ are empty.
\begin{center}
\begin{tikzpicture}[x=0.82cm, y=0.82cm]
    \tikzset{gv/.append style={minimum size=6.2mm, font=\scriptsize\bfseries}}
    \begin{scope}[g2p5-allrev, g2p5-a6/.style=gvcur, g2p5-a3/.style=gvseen, g2p5-a4/.style=gvseen,
                  g2p5-f36/.style=g2p5-rt, g2p5-f46/.style=g2p5-rt]
        \pic {g2p5-S1};
        \node[font=\scriptsize\bfseries] at (1.6,2.8) {1. pop 6: mark 3, 4};
    \end{scope}
    \begin{scope}[xshift=4.0cm, g2p5-allrev, g2p5-a6/.style=gvdst, g2p5-a3/.style=gvseen, g2p5-a4/.style=gvcur,
                  g2p5-a0/.style=gvseen, g2p5-a2/.style=gvseen,
                  g2p5-f36/.style=g2p5-rt, g2p5-f46/.style=g2p5-rt, g2p5-f04/.style=g2p5-rt, g2p5-f24/.style=g2p5-rt]
        \pic {g2p5-S1};
        \node[font=\scriptsize\bfseries] at (1.6,2.8) {2. pop 4: mark 0, 2};
    \end{scope}
    \begin{scope}[xshift=8.0cm, g2p5-allrev, g2p5-a6/.style=gvdst, g2p5-a3/.style=gvcur, g2p5-a4/.style=gvdone,
                  g2p5-a0/.style=gvdone, g2p5-a2/.style=gvdone, g2p5-a1/.style=gvseen,
                  g2p5-f36/.style=g2p5-rt, g2p5-f46/.style=g2p5-rt, g2p5-f04/.style=g2p5-rt, g2p5-f24/.style=g2p5-rt,
                  g2p5-f13/.style=g2p5-rt]
        \pic {g2p5-S1};
        \node[font=\scriptsize\bfseries] at (1.6,2.8) {3. pops 2, 0; pop 3: mark 1};
    \end{scope}
    \begin{scope}[xshift=12.0cm, g2p5-allrev, g2p5-a7/.style=gvdst, g2p5-a3/.style=gvdone, g2p5-a2/.style=gvdone,
                  g2p5-a0/.style=gvdone, g2p5-a1/.style=gvdone,
                  g2p5-f37/.style=g2p5-rt, g2p5-f27/.style=g2p5-rt, g2p5-f03/.style=g2p5-rt, g2p5-f13/.style=g2p5-rt]
        \pic {g2p5-S1};
        \node[font=\scriptsize\bfseries] at (1.6,2.8) {4. the search from 7};
    \end{scope}
\end{tikzpicture}
\end{center}
\begin{center}\small
\begin{tabular}{@{}cllll@{}}
\toprule
\textbf{search} & \textbf{pop} & \textbf{newly marked} & \textbf{stack after} & \textbf{note} \\
\midrule
$v = 6$ & 6 & 3, 4 & \texttt{3 4} & \\
        & 4 & 0, 2 & \texttt{3 0 2} & \\
        & 2 & --- & \texttt{3 0} & in-degree 0 \\
        & 0 & --- & \texttt{3} & \\
        & 3 & 1 & \texttt{1} & 0 is already marked \\
        & 1 & --- & empty & scan: \textbf{0 1 2 3 4} \\
\midrule
$v = 7$ & 7 & 2, 3 & \texttt{2 3} & \\
        & 3 & 0, 1 & \texttt{2 0 1} & \\
        & 1, 0, 2 & --- & empty & scan: \textbf{0 1 2 3} \\
\midrule
$v = 5$ & 5; 3; 1, 0 & 3; 0, 1; --- & & scan: \textbf{0 1 3} \\
\bottomrule
\end{tabular}
\end{center}
The visiting order from 6 is $3, 4, 0, 2, 1$; the index-order scan turns it into the sorted list. Over all eight searches $0+0+0+2+2+3+5+4 = 16$ ancestors are written --- exactly the size of the official output.
\end{tracebox}

\subsubsection*{6. Test suite}

\begin{testbox}
\centering\small
\begin{tabular}{@{}p{5.2cm}p{4.2cm}p{4.4cm}@{}}
\toprule
\textbf{Input} & \textbf{Our output} & \textbf{Why} \\
\midrule
$n = 1$, no edges & \texttt{[[]]} & smallest input \\
$n = 3$, no edges & \texttt{[[],[],[]]} & every list empty \\
chain $0 \to 1 \to 2 \to 3$ & \texttt{[[],[0],[0,1],[0,1,2]]} & long reverse paths \\
chain $3 \to 2 \to 1 \to 0$ & \texttt{[[1,2,3],[2,3],[3],[]]} & visiting order $\ne$ sorted order \\
diamond $0{\to}1, 0{\to}2, 1{\to}3, 2{\to}3$ & \texttt{[[],[0],[0],[0,1,2]]} & 0 reaches 3 twice, listed once \\
$n = 6$: $0{\to}1$, $2{\to}3$, $4{\to}5$ & \texttt{[[],[0],[],[2],[],[4]]} & disconnected \\
chain of 1000 $+$ 1001 random edges & 499\,500 entries, 13\,ms & the largest output \\
\bottomrule
\end{tabular}
\end{testbox}
All outputs come from running Approaches 2 and 3 (they agree). In addition, Approaches 1--4 were run on 3000 random DAGs with $n \le 9$ (a random vertex order, each forward pair an edge with probability $1/3$, edges shuffled) and produced identical lists every time.

\subsubsection*{7. Common mistakes}

\begin{warningbox}[Mistake 1 --- forward search without per-source marks]
Appending $u$ on every \emph{path} instead of every \emph{vertex} gives duplicates: without \texttt{mark} the forward version prints \texttt{[0,0,1,2,3,4]} for vertex 6 of Sample~1 (0 reaches 6 via 3 and via 4). Worse, the running time becomes the number of paths: a ladder of $k$ layers of two vertices, each layer joined completely to the next, has $2^{k-1}$ paths from a top vertex to a bottom vertex, and $m \le 2000$ allows $k = 500$.
\end{warningbox}

\begin{warningbox}[Mistake 2 --- collecting in visiting order]
Pushing $u$ into the answer when it is \emph{marked} produces \texttt{[3,4,0,2,1]} for vertex 6. Either scan the marks in index order (Approach~2), or let the outer loop run in increasing order (Approach~3), or sort.
\end{warningbox}

\begin{warningbox}[Mistake 3 --- a fresh \texttt{vector<bool> seen(n)} per search, inside a recursive helper that copies it]
A new array per search is fine (\bigO{n} each, \bigO{n^2} total). Passing it \emph{by value} into a recursive DFS copies it at every call: \bigO{n} per visited vertex, up to \bigO{n^3} $= 10^9$ in total, and \tle. Pass by reference, or use the \texttt{mark[u] = v} trick.
\end{warningbox}

\begin{warningbox}[Mistake 4 --- including $v$ itself]
The search from $v$ marks $v$. ``Ancestor'' requires $u \ne v$: skip it in the scan.
\end{warningbox}

\subsubsection*{8. Takeaways}

\begin{itemize}[leftmargin=*]
    \item ``Everything that reaches $v$'' is ``everything reachable from $v$ in $G^R$''.
    \item When the output is quadratic, $n$ linear searches are optimal up to a constant.
    \item Reuse one marker array across searches with \texttt{mark[u] = v}.
    \item Sorted output comes for free from index-order scans or an increasing outer loop.
\end{itemize}

\begin{variantsbox}
\variant{Only the counts}{\varnone}
  {return, for every $v$, the number of ancestors instead of the list.}
  {the same reachability sets; only their sizes are needed.}
  {count in the same $n$ reverse searches. (Summing ancestor counts of a DAG by DP double-counts shared ancestors --- on the diamond it gives 4 for vertex 3 instead of 3.)}
\variant{Descendants instead}{\varnone}
  {for every $v$, the vertices $v$ can reach.}
  {no reversal: forward searches in $G$.}
  {Approach~2 with \texttt{adj} in place of \texttt{radj}; or Approach~3 with the roles of the two lists swapped.}
\variant{Ancestors of one vertex}{\varnone}
  {only one query vertex $q$.}
  {a single reverse search.}
  {template 1 of the pattern with $T = \{q\}$: \VE.}
\variant{Cycles allowed}{\varnone}
  {the graph is a general directed graph.}
  {``ancestor'' still means ``can reach''; vertices on a cycle through $v$ become ancestors of $v$.}
  {Approaches 1--3 are unchanged (no acyclicity was used); Approach~4 breaks, since Kahn's order misses the cycle vertices --- it would first need the strongly connected components (Chapter~19).}
\variant{Common ancestors}{\varnone}
  {for a pair $(a, b)$, the vertices that reach both.}
  {two target sets.}
  {two reverse searches, from $a$ and from $b$, and an AND of the marks (like E5.2).}
\end{variantsbox}
